\section{Introduction}\label{sec:intro}

Estimating a linear function of the state, rather than the whole state, is often what a
control loop actually needs. Functional observers do exactly this, and their appeal is
quantitative: the order of a functional observer can be far below the order of a
Luenberger observer for the same system, because only the directions that the functional
sees have to be estimated. The design question is therefore an order question: given the
system and the functionals to be estimated, what is the smallest number of observer rows
for which a functional observer exists?

For a linear system with measured output, Fernando, Trinh and Jennings \cite{fer10}
reduced that question to two rank conditions on the observer matrix. Writing the observer
rows as a block $L$ that contains the given functional rows $L_0$, the conditions are a
rank equality for the algebraic part and a rank equality for every complex value of the
spectral variable. They characterize the feasible rows for a prescribed order, but they do
not address the cost of deciding whether such rows exist.

The characterization itself has since been completed. A unified algebraic framework is
given in \cite{dar25}. The two recent works \cite{fer26a,fer26b} determine, order by order,
which orders are attainable: \cite{fer26a} identifies through functional observability
indices the minimum dimension permitted by the algebraic condition, the complete family of
augmentations attaining it, and necessary and sufficient spectral conditions for that
dimension to be optimal, while \cite{fer26b} characterizes the interval between that
minimum and the endpoint set by functional observability, using a nullspace representation
of the algebraic condition and a matrix-pencil form of the spectral one. Neither work
states a decision procedure with a running-time or bit-complexity bound, and the per-order
conditions are existential statements over a chart and a parameter. The complexity of
deciding feasibility at a prescribed order is therefore left open by the characterization
literature. This paper takes the conditions as given and studies that complexity.

Adjacent design problems are known to be hard, which makes the question natural but does
not answer it. Choosing a minimal set of sensors that renders a system functionally
observable, or structurally functionally observable, is NP-hard \cite{zha25}, and
constrained sensor-placement variants are NP-complete \cite{dey21}. Static output feedback
with a spectral predicate is NP-hard \cite{lof26,bar26} and, in its real feasibility form,
$\exists\mathbb R$-complete \cite{liu26}; affine rank minimization is
$\exists\mathbb R$-complete already at low rank \cite{maj26}. None of these results
transfers. In sensor placement the quantified object is a subset of the rows of the output
matrix, so the output matrix is itself a design variable, whereas here it is fixed and the
design variable is the set of observer rows; the predicates also differ, a condition on the
spectral abscissa of a closed loop against a constant-rank condition that has to hold for
every complex value of the spectral variable. A reduction would moreover have to preserve
the single block of free rows that occurs in both conditions, together with the common
origin of the compressed matrices, whereas rank-constrained formulations leave every entry
of the constrained matrix free.

The approach of this paper is to make the decision problem finite-dimensional. Restricting
row functionals to the common kernel of the output matrix and its product with the system
matrix removes exactly the directions that those two matrices already account for, and
turns the two conditions into a row-space inclusion and into the constant rank of a
pencil. Every free observer row is then determined by its coordinates in the space of
attainable row pairs, whose dimension is the effective dimension
$\kappa_{\rm obs}=r\operatorname{rank}[N\ AN]$ of the free data, with $r$ the number of
free rows and $N$ a basis of the common kernel. Two consequences are developed. First, the
spectral condition is equivalent to the non-vanishing of a finite algebraic certificate,
and Cauchy--Binet compresses that certificate to a matrix whose size depends on
$\kappa_{\rm obs}$ rather than on the number of maximal minors. Second, deciding
feasibility becomes a real-algebraic feasibility problem whose number of variables is
controlled by either the prescribed order or the effective dimension, which yields the two
parameterized upper bounds of this paper. The boundary case of no free rows is decided
separately, by a finite rational linear system obtained from a polynomial left inverse.

\subsection{Contributions and organization}\label{sec:contrib}

The contributions are the following.

\begin{enumerate}
\item A compression of the two observer conditions to a row-space inclusion and to the
constant rank of a pencil, together with the identification of the effective dimension
$\kappa_{\rm obs}$ of the free observer data as the structural parameter of the problem.
\item A finite algebraic certificate for the spectral condition, and its compression by
Cauchy--Binet to a matrix of size controlled by $\kappa_{\rm obs}$; no enumeration of
maximal minors is required.
\item A decision algorithm for every fixed order, with a bit-complexity bound
$\mathcal L^{O(q^2)}$ in the input length $\mathcal L$ and the prescribed order $q$.
\item A decision algorithm that is fixed-parameter tractable in $\kappa_{\rm obs}$, with
bound $f(\kappa_{\rm obs})\operatorname{poly}(\mathcal L)$, together with an exact
polynomial-time test for the zero-free-row case.
\item Membership of the prescribed-order decision problem, and of its finite threshold
version, in the existential theory of the reals.
\end{enumerate}

The paper is organized as follows. Section~\ref{sec:prelim} fixes the notation and states
the decision, optimization and construction problems. Section~\ref{sec:compressed} proves
the compression, Section~\ref{sec:dimension} identifies the effective dimension, and
Section~\ref{sec:certificate} develops the certificate. Sections~\ref{sec:fixedq}
and~\ref{sec:fpt} give the two algorithms, Section~\ref{sec:existsr} records the
existential-real membership, and Section~\ref{sec:conclusions} concludes.

\section{Preliminaries}\label{sec:prelim}

\subsection{Notation and conventions}\label{sec:notation}

The system is $(A,C)$ with $A\in\mathbb Q^{n\times n}$ and $C\in\mathbb Q^{p\times n}$;
the given functional rows form $L_0\in\mathbb Q^{r_0\times n}$, and $q\ge r_0$ is the
prescribed observer order. The observer rows are the rows of $L=[L_0;R]$, where
$R\in\mathbb R^{r\times n}$ collects the $r=q-r_0$ free rows; throughout, $L_0$ is fixed
data and only $R$ is free. All input data are rational and encoded in binary, and
$\mathcal L$ denotes the total input bit length; $q\le n$ throughout. A matrix $M$ acts on
row vectors by $x\mapsto xM$; $\operatorname{rank}M$ is the dimension of its row space and
$\operatorname{row}M$ that row space; $\ker M$ is the kernel of the map $v\mapsto Mv$. For
a matrix pencil $P(\lambda)$, its $\omega\times\omega$ minors are the maximal minors when
$\operatorname{rank}P(\lambda)=\omega$. Ranks are computed over $\mathbb C$, which for real
matrices agrees with the rank over $\mathbb R$; the spectral variable is complex. If a
complex observer is allowed, realification doubles the free coordinates and replaces every
transpose by the conjugate transpose, which changes no bound below.

\subsection{The observer conditions}\label{sec:conditions}

A functional observer of order $q$ is a matrix $L=[L_0;R]\in\mathbb R^{q\times n}$ such
that
\begin{equation}\label{eq:5a}
\operatorname{rank}\begin{bmatrix}LA\\ L\\ CA\\ C\end{bmatrix}
=\operatorname{rank}\begin{bmatrix}L\\ CA\\ C\end{bmatrix},
\end{equation}
\begin{equation}\label{eq:5b}
\operatorname{rank}\begin{bmatrix}\lambda L-LA\\ CA\\ C\end{bmatrix}
=\operatorname{rank}\begin{bmatrix}L\\ CA\\ C\end{bmatrix}
\qquad\text{for every }\lambda\in\mathbb C.
\end{equation}
These are the Fernando--Trinh--Jennings conditions \cite{fer10} (conditions (5a) and (5b)
in the notation of that paper). Condition \eqref{eq:5a} is the algebraic one and
\eqref{eq:5b} is the spectral one; the latter has to hold at every complex value of
$\lambda$, including the finite spectrum of $A$.

\subsection{Decision, optimization and construction}\label{sec:problems}

The three questions associated with \eqref{eq:5a}--\eqref{eq:5b} are different, and this
paper studies the first. The decision problem $\mathrm{DEC}$ asks, for given
$(A,C,L_0)$ and $q$, whether some $R$ satisfies \eqref{eq:5a} and \eqref{eq:5b}; a witness
is not required. The optimization problem $\mathrm{OPT}$ asks for the least such $q$, whose
value is written $q_{\min}$, and the construction problem $\mathrm{FIND}$ asks for a
witness $R$ at that order. Throughout, $q1$ denotes the complexity of the \emph{decision}
problem $\mathrm{DEC}$. The threshold problem
\begin{equation}\label{eq:mindec}
\mathrm{MINDEC}(Q)=\bigvee_{q=r_0}^{\min(Q,n)}\mathrm{DEC}(q)
\end{equation}
asks whether some order up to $Q$ is feasible; it is a finite disjunction of decision
instances. The algorithms below decide $\mathrm{DEC}$, and when they answer affirmatively
the algebraic certificate they construct contains a witness, but no claim is made here
about the complexity of $\mathrm{OPT}$ or $\mathrm{FIND}$, and monotonicity of
$\mathrm{DEC}$ in $q$ is not used.
